% =====================================================================
%  Project Introduction  --  Gezhi
%  2026 AI for SDGs: Global Youth AI Future Innovation Competition
%  Macau Secondary School Division
%
%  Compile:  pdflatex ProjectIntroduction_Gezhi.tex   (run twice)
%  English-only, as required by the competition rules.
% =====================================================================
\documentclass[10pt,a4paper]{article}

\usepackage[T1]{fontenc}
\usepackage[utf8]{inputenc}
\usepackage{lmodern}
\usepackage[margin=1.9cm,top=1.7cm,bottom=1.7cm]{geometry}
\usepackage{booktabs}
\usepackage{tabularx}
\usepackage{array}
\usepackage{enumitem}
\usepackage{titlesec}
\usepackage{xcolor}
\usepackage{microtype}
\usepackage[hidelinks]{hyperref}

% ---- compact section headings -------------------------------------
\titleformat{\section}{\normalfont\large\bfseries\color{black!85}}{\thesection.}{0.5em}{}
\titlespacing*{\section}{0pt}{7pt}{3pt}
\titleformat{\subsection}[runin]{\normalfont\bfseries}{}{0em}{}[\hspace{0.5em}]
\titlespacing*{\subsection}{0pt}{5pt}{0.5em}

\setlist[itemize]{leftmargin=1.1em,itemsep=1pt,topsep=2pt,parsep=0pt}
\setlength{\parskip}{2pt}
\setlength{\parindent}{0pt}
\pagestyle{empty}

\newcommand{\gezhi}{\textsc{Gezhi}}
\newcommand{\xiaoge}{\textsc{Xiaoge}}
\newcommand{\fillin}[1]{\textcolor{red}{\textbf{[TO BE COMPLETED:\,#1]}}}

\begin{document}

% ===================== TITLE BLOCK ==================================
\begin{center}
{\LARGE\bfseries Gezhi: A Local-First AI Study Companion\\[2pt]
That Puts the Student in the Teacher's Chair}\\[6pt]
{\small 2026 AI for SDGs --- Global Youth AI Future Innovation Competition\\
Macau Secondary School Division \quad$\vert$\quad Theme: AI and Education}\\[4pt]
{\footnotesize Team \textit{[Team Name]} $\cdot$ \textit{[School Name]} $\cdot$ Macau}
\end{center}
\vspace{-2pt}
\hrule
\vspace{4pt}

% ===================== 1. PROBLEM ===================================
\section{The Problem}

Generative AI has been promoted as the technology that will finally give every student a
personal tutor. The first rigorous evidence says otherwise, and it points in two opposite
directions at once.

\begin{table}[h]
\centering
\footnotesize
\begin{tabularx}{\textwidth}{@{}p{2.9cm} p{3.4cm} X@{}}
\toprule
\textbf{Study} & \textbf{Design} & \textbf{Finding} \\
\midrule
Bastani et al.\ (2024) &
Randomised controlled trial; $\sim$1{,}000 Turkish high-school students; mathematics &
With unrestricted ChatGPT: \textbf{+48\%} correct on practice problems, but \textbf{\(-17\%\)}
on the exam once the AI was removed. With a tutor engineered to give hints but never
answers: \textbf{+127\%} on practice, yet \textbf{no exam gain} over the no-AI control. \\[2pt]
Oreopoulos \& Low (2026) &
Two-year cluster RCT; 18 Tennessee middle schools; AI configured to coach, not answer &
Effect of \textbf{0.06--0.08 SD} --- comparable to using the same platform \emph{without}
AI. The binding constraint was behaviour: the median student messaged the tutor on only
\textbf{one third} of practice days, and in only \textbf{17\%} of sessions in which they
made a mistake. \\
\bottomrule
\end{tabularx}
\end{table}

\vspace{-4pt}
Read together, these results describe a narrow gap. \textbf{AI tutors today either harm
learning or go unused.} Making the AI more careful does not fix the second problem: even a
deliberately restrained tutor produced no measurable benefit, because students did not
open it. The only condition that has never been shown to harm learning is the one in which
\textbf{the student produces the explanation}.

% ===================== 2. USERS =====================================
\section{Target Users}

\textbf{Primary:} secondary-school students in Macau (Secondary 3--6, roughly ages 14--18)
preparing for school examinations and university entrance. Our specific target is the
student who \emph{feels} they understand a topic but cannot explain it. This ``illusion of
knowing'' is not anecdotal: across 115 studies and 15{,}889 learners, the correlation
between students' own judgements of what they have learned and their actual performance is
only \(r = .178\) (Yang et al., 2022).

\textbf{Secondary:} subject teachers who need evidence of what a class genuinely
understands rather than what it has completed; and parents who want something more
informative than a study-time counter.

% ===================== 3. FEATURES ==================================
\section{Main Features}

\gezhi{} is a Windows desktop study companion organised around one loop:
\textbf{expose the gap, explain it, defend it, master it.}

\begin{itemize}
  \item \textbf{Knowledge vault.} The student's own course materials --- textbooks,
  handouts, notes, scanned worksheets --- indexed locally and searchable in both Chinese
  and English. Nothing is uploaded.
  \item \textbf{Question bank and exam builder.} Questions tagged by knowledge point and
  exportable as a properly typeset examination paper, with an \emph{interleaved} mode that
  mixes topics and deliberately withholds topic labels.
  \item \textbf{Explain-and-defend (the core).} The student explains a topic aloud to
  \xiaoge{}, an on-screen companion who is deliberately \emph{a year below them} and who
  holds a fixed set of misconceptions. \xiaoge{} does not nod along --- it argues back. The
  student is assessed on a six-dimension rubric and on how many rounds it takes to actually
  convince it.
  \item \textbf{Spaced review and planning.} Every knowledge point carries a mastery
  estimate that schedules both questions \emph{and} explanations.
  \item \textbf{Progress without scores.} \xiaoge{} carries a three-by-three lattice above
  its head. Nine empty cells mean it understands nothing yet; each successful explanation
  lights one cell. Students see progress without reading a number.
\end{itemize}

% ===================== 4. AI METHODS ================================
\section{How the AI Works}

\gezhi{} does not implement its own agent loop. The reasoning layer is \textbf{DeepSeek
Harness (DSH)}, an MIT-licensed agent kernel invoked as a controlled subprocess that emits
a newline-delimited JSON event stream. This provides session persistence, tool sandboxing,
context compaction and an approval policy that a student team could not responsibly rebuild.
Three decisions matter most:

\textbf{1. The vault is the sandbox.} The subprocess working directory is set to the
student's vault, and the kernel derives its write boundary from that directory. The AI can
read and write the student's own notes and nothing else --- a security boundary enforced by
the operating system's process model rather than by prompt instructions.

\textbf{2. Chinese retrieval required a non-obvious fix.} We measured that SQLite's default
FTS5 tokeniser returns \emph{zero} results for Chinese substring queries, because it treats
an entire Chinese sentence as a single token. Switching to \texttt{tokenize='trigram'}
yields correct matches with 1\,ms query latency over 200 documents. Retrieval is therefore
local, offline and citation-bearing: every claim the system makes is anchored to a page of
the student's own textbook.

\textbf{3. The companion is engineered to be wrong on purpose.} A misconception library,
derived from the student's own error statistics, is injected into \xiaoge{}'s system prompt
together with a hard constraint: it may not concede until the student's explanation actually
addresses the misconception. This follows CHI 2025 work showing that a teachable agent must
\emph{play dumb} to be useful, and a controlled study finding that when an AI simply
supplied correct answers, students' error-correction ability did not improve at all.

% ===================== 5. NOVELTY ===================================
\section{What Is Novel}

We do not claim to have invented learning by teaching. Teachable agents date back to Chase
et al.\ (2009), and 2024--2026 has seen active research in this area. Our contribution is
four specific things:

\begin{enumerate}[leftmargin=1.4em,itemsep=1pt,topsep=2pt]
  \item That research has overwhelmingly involved \textbf{university} students, largely in
  computer science. We bring it to \textbf{secondary-school} learners, in a
  Chinese-language desktop product built for their own curriculum.
  \item The student teaches \textbf{the knowledge points in their own textbook}, retrieved
  locally with page-level citations --- not generic exercises.
  \item We convert ``convincing \xiaoge{}'' into a \textbf{measurable score that feeds a
  spaced-repetition schedule}, closing a loop rather than producing a one-off conversation.
  \item The knowledge base and learning record \textbf{never leave the machine}; model calls
  run through the user's own account and prompts are not retained.
\end{enumerate}

% ===================== 6. IMPACT ====================================
\section{Expected Impact}

We measure success on learning outcomes, not engagement. Three instruments: (a) the change
in Feynman rubric score for a knowledge point before and after a teach-and-defend cycle;
(b) the \textbf{calibration gap} between a student's self-rating and their rubric score,
tracked over time; (c) the number of dialogue turns required to resolve a misconception,
which should fall as understanding deepens.

\gezhi{} addresses \textbf{SDG 4 (Quality Education)} by making explanatory practice ---
normally available only through one-to-one tutoring --- accessible to any student with a
laptop, and by treating \textbf{AI literacy} as a learning objective in its own right:
students who diagnose and correct an AI's misconceptions learn first-hand where such systems
fail. This maps directly onto the \emph{AI system design} dimension at the \emph{Create}
level of UNESCO's AI Competency Framework for Students (2024).

% ===================== 7. STATUS ====================================
\section{Current Status}

A working prototype runs on Windows and demonstrates the full loop described above:
local indexing of a mixed Chinese-and-English material set, citation-bearing retrieval, a
spoken explanation session with the \xiaoge{} companion, rubric scoring, and export of a
typeset examination paper. The misconception library currently holds \fillin{number}
entries across \fillin{subject areas}. A first round of use with \fillin{number}
secondary-school students has been completed, and the measured outcomes are reported in the
accompanying Research Report.

Our immediate next steps are to add a comparison condition without the AI, to validate the
scoring rubric against human raters, and to run the study over a period long enough for the
spaced-repetition schedule to matter.

\end{document}
